\section{Bodies, Place, and Motion}\label{sec:bodies}

Having treated the angelic substance in itself, Aquinas considers it in
comparison with corporeal things in three questions: the angels compared
with bodies, with corporeal places, and with local motion (\latin{primo, de
comparatione Angelorum ad corpora; secundo, de comparatione Angelorum ad
loca corporalia; tertio, de comparatione Angelorum ad motum localem};
\ST{I q.~51 pr.}). The three questions apply one principle drawn from
q.~50 (\S\ref{sec:q50}). A pure intellect has no dimensive quantity, so
everything bodily it has or does comes from its power acting on a body,
not from any body belonging to its nature. On that principle Aquinas
explains how Scripture's angels can be seen, eat, walk, stand in a place,
and travel, while the angel remains a spirit.

\subsection{The Angels in Comparison with Bodies (\ST{I q.~51})}\label{sec:q51}

The question has three articles: whether the angels have bodies naturally
united to them, whether they assume bodies, and whether they exercise
functions of life in the bodies they assume (\ST{I q.~51 pr.}).

\paragraph{No body by nature (a.~1).} The \emph{sed contra} is Dionysius,
\emph{Divine Names}~4, that ``the angels are understood to be
incorporeal.'' Aquinas reasons that what belongs to a nature only
accidentally is not found in every instance of it. Since understanding is
not the act of a body, having a body is not of the nature of intellectual
substance as such. It belongs to the human soul because the soul is
imperfect, ``not having the fulness of knowledge in its own nature, but
acquiring it from sensible things through the bodily senses.'' And
``whenever we find something imperfect in any genus we must presuppose
something perfect in that genus''; therefore ``some are quite separated
from bodies, and these we call angels'' (\latin{aliquae sunt a corporibus
separatae. Et has dicimus Angelos}). The objections gather four
authorities. Origen's own sentence, in Rufinus's Latin as the
\emph{Ante-Nicene Fathers} translate it, is that ``it is an attribute
of the divine nature alone---i.e., of the Father, Son, and Holy
Spirit---to exist without any material substance, and without partaking in
any degree of a bodily adjunct'' (\emph{De principiis} I.6.4). Aquinas
judges that Origen followed the philosophers who admitted no incorporeal
substance except as united to a body, ``being deceived here as he was also
in many other points'' (ad~1). Bernard's saying that every created spirit
needs a body is taken of a bodily instrument assumed for a purpose.
Augustine's ``animals of the atmosphere'' (\emph{De Genesi ad litteram}
III, as the \emph{Summa} cites it) reports the Platonists' opinion without
asserting it (ad~1). Augustine elsewhere leaves that opinion undecided:
``I dare not determine whether there be some spirits embodied in an aerial
substance'' (\emph{De civitate Dei} XV.23). Gregory's ``rational animal''
(\emph{Hom.\ in Ev.}\ 10.1) is a metaphor (ad~2). To give life to a body
formally belongs to a substance that is only part of a nature, so ``an
intellectual substance which is not united to a body is more perfect than
one which is united to a body'' (ad~3).

\paragraph{Assumed bodies (a.~2).} The \emph{sed contra} is Augustine,
\emph{De civitate Dei}~XVI, that angels appeared to Abraham in assumed
bodies (\latin{in assumptis corporibus}). Augustine's chapter opens: ``God
appeared again to Abraham at the oak of Mamre in three men, who it is not
to be doubted were angels'' (\emph{De civ.\ Dei} XVI.29). He also affirms
that ``angels have appeared to men in such bodies as could not only be
seen, but also touched'' (XV.23). Aquinas rejects the opinion that every
angelic apparition in Scripture happened ``in prophetic vision---that is,
according to imagination.'' What is seen in imaginary vision is in the
beholder alone, whereas Scripture presents angels ``seen commonly by all'':
by Abraham ``and by his whole family, by Lot, and by the citizens of
Sodom,'' and by all who were with Tobias. Such sight is bodily vision, which
only a body can terminate; the angels are not bodies and have none by
nature; ``it follows that they sometimes assume bodies.'' The replies
define the \latin{corpus assumptum}. Angels need it ``not for themselves,
but on our account; that by conversing familiarly with men they may give
evidence of that intellectual companionship which men expect to have with
them in the life to come'' (\latin{intelligibilem societatem}). The
apparitions of the Old Law were also ``a figurative indication that the
Word of God would take a human body'' (ad~1). The body is united to the
angel ``not as its form, nor merely as its mover, but as its mover
represented by the assumed movable body.'' As Scripture sets forth
intelligible things under sensible likenesses, so ``sensible bodies are so
fashioned by angels as fittingly to represent the intelligible properties
of an angel'' (ad~2). The matter is air, condensed ``by the Divine power in
so far as is needful for forming the assumed body'' (ad~3).

\paragraph{No functions of life in assumed bodies (a.~3).} The objections
are drawn from Scripture's own scenes. Abraham ``walked with them, bringing
them on the way'' (Gen 18:16). Raphael told Tobias, ``I have often walked
through all the ways thereof'' (Tob 5:8). The angels spoke, and at Mamre
``they had eaten'' (Gen 18:9). The \emph{sed contra} answers from a.~2:
bodies assumed by angels do not live. Aquinas determines that vital acts
share something with non-vital ones, as speech is a sound and walking a
motion. Angels can perform in assumed bodies what is common to both, ``but
not as to that which is special to living subjects,'' because ``that which
has the faculty has the action'' (\latin{cuius est potentia, eius est
actio}). The replies apply this to each scene. There is no pretence, any
more than when Scripture figures spiritual things by bodily ones. The
bodies are assumed ``that the spiritual properties and works of the angels
may be manifested by the properties of man and of his works'' (ad~1). The
angel perceives nothing through the organs; the organs signify its powers,
``just as by the eye the power of the angel's knowledge is pointed out
\dots\ as Dionysius teaches'' (ad~2). Dionysius writes of the angelic
figures that ``the powers of vision denote the most transparent elevation
towards the Divine lights'' (\emph{CH}~15.3). The angel is moved only
accidentally with the body it moves (ad~3), and speaks only by ``a
semblance of speech,'' forming sounds in the air (ad~4). Eating is
``not a true eating, but figurative of spiritual eating,'' because the food
was not converted into the assumed body. Raphael says so: ``I seemed indeed
to eat and to drink with you but I use an invisible meat and drink, which
cannot be seen by men'' (Tob 12:19, Douay). Abraham served the guests as
men and worshipped God in them (ad~5). Augustine says the same: Abraham's household ``could not doubt that God was in them as He was
wont to be in the prophets'' (\emph{De civ.\ Dei} XVI.29). On Gen 6:4
Aquinas follows Augustine, who identifies the ``sons of God, who are also
called angels of God'' with ``the sons of Seth'' and the daughters of men
with ``the daughters of Cain'' (\emph{De civ.\ Dei} XV.23). He adds that
if any are begotten through demons, it is neither from a demon's seed nor
from its assumed body, ``so that the person born is not the child of a
demon, but of a man'' (ad~6).

\angeldossier{\ST{I q.~51 aa.~1--3}.}
 {Church teaching: the angels are ``purely spiritual creatures''
 (CCC~330), and so have no body by nature (a.~1). Common teaching: angels
 have appeared in bodies seen and touched by many, not only in imaginary
 vision (a.~2; Gen 18--19; Tob~5, 12; Augustine, \emph{De civ.\ Dei}
 XV.23, XVI.29), and they did not truly eat (a.~3 ad~5; Tob 12:19).
 Thomist position: the assumed body is of condensed air, united to the
 angel as a represented mover, and performs no vital act (a.~2 ad~2--3;
 a.~3).}
 {Dionysius, \emph{DN}~4 (\emph{sed contra} of a.~1) and \emph{CH}~15.3
 (a.~3 ad~2); Augustine, \emph{De civ.\ Dei} XV.23 and XVI.29 (\emph{sed
 contra} of a.~2; a.~3 ad~5--6), \emph{De Gen.\ ad litt.} III and
 \emph{De Trin.} (as the \emph{Summa} cites them);
 Origen, \emph{De principiis} I.6.4; Bernard, \emph{In Cant.}\ serm.~6 and
 Gregory, \emph{Hom.\ in Ev.}\ 10.1; Gen 6:4, 18,
 19; Tob 5:7--8, 12:19; Luke~24; 1~Kgs [1~Sam] 2:6; CCC~330.}
 {Origen: incorporeity belongs to the Trinity alone (a.~1); the Platonist
 opinion of aerial animals, reported by Augustine (a.~1 ad~1); those who
 held that all angelic apparitions were imaginary visions (a.~2); the
 reading of Gen 6:4 as angels' unions with women, which Augustine rejects
 (a.~3 ad~6).}

\subsection{The Place of the Angels (\ST{I q.~52})}\label{sec:q52}

Three articles ask whether an angel is in a place, whether it can be in
several places at once, and whether several angels can be in the same
place (\ST{I q.~52 pr.}).

\paragraph{In place by power (a.~1).} The objections deny that an angel is
in place at all, citing Boethius and Aristotle that only bodies are. The
\emph{sed contra} is liturgical: the collect of Compline, ``Let Thy holy
angels who dwell herein, keep us in peace'' (\latin{Angeli tui sancti,
habitantes in ea, nos in pace custodiant}). Aquinas determines that ``it is
befitting an angel to be in a place; yet an angel and a body are said to be
in a place in quite a different sense.'' A body is in place by contact of
dimensive quantity; the angel has none, ``for theirs is a virtual one''
(\latin{sed est in eis quantitas virtualis}). ``Consequently an angel is
said to be in a corporeal place by application of the angelic power in any
manner whatever to any place'' (\latin{per applicationem virtutis angelicae
ad aliquem locum qualitercumque}). The angel is therefore neither
commensurate with its place nor contained by it. It contains its place,
``for the soul is in the body as containing it, not as contained by it.''
The \latin{locus} of an angel is whatever body its power touches, and the
objections, which suppose place in the bodily sense, fall away.

\paragraph{Not in several places at once (a.~2).} Damascene gives the
\emph{sed contra}: ``They are circumscribed: for
when they are in the Heaven they are not on the earth: and when they are
sent by God down to the earth they do not remain in the Heaven.'' He adds
that, wherever they are assigned, ``there they are present after the
manner of a mind and energise, and cannot be present and energise in
various places at the same time'' (\emph{De fide orthodoxa} II.3). Aquinas
gives the reason. God's power is infinite and touches all things, but
``since the angel's power is finite, it does not extend to all things, but
to one determined thing.'' Since the angel is in place by applying that power, it
is ``not everywhere, nor in several places, but in only one place.'' This
one place need not be a point, as some thought who could not get beyond
imagination, since the angel is ``beyond the genus of quantity and
situation.'' Its place may be great or small ``according as to his own
free-will he applies his power to a great or to a small body.'' Aquinas
then fixes the three modes of presence: a body is in place
\latin{circumscriptive}, because it is measured by the place; the angel
\latin{definitive}, ``because he is in a place in such a manner that he is
not in another''; God neither, ``because He is everywhere.'' The objections
from the soul whole in every part of the body (Augustine, \emph{De
Trinitate} VI) and from the angel's several effects at Sodom (Gen 19:25)
are answered together: ``the entire subject to which the angelic power is
immediately applied, is reputed as one place, even though it be
continuous.'' Augustine himself says of the soul that ``in each
body, it is both whole in the whole, and whole in each several part of
it'' (\emph{De Trin.} VI.6.8).

\paragraph{Not two angels in one place (a.~3).} The \emph{sed contra} runs
by analogy: there are not two souls in one body. Aquinas's reason is causal
rather than spatial. Angels do not fill a place, so crowding is not the
bar (ad~1). The bar is that ``it is impossible for two complete causes to
be the causes immediately of one and the same thing'' (\latin{impossibile
est quod duae causae completae sint immediatae unius et eiusdem rei}).
Several men rowing a boat are no exception, because none of them is a
perfect mover and together they act as one. Since the angel is in a place
by its power touching that place immediately ``by way of a perfect
container,'' there can be but one angel in one place. A demon and a soul
in the same body are no counter-instance, since they are not related to it
by the same causality: the soul is its form and the demon is not (ad~3).

\angeldossier{\ST{I q.~52 aa.~1--3}.}
 {Common teaching: an angel is in place not by dimension but by
 operation, and not in several places at once (aa.~1--2; Damascene,
 \emph{De fide orth.} II.3). Thomist position: presence by
 \latin{applicatio virtutis} and \latin{quantitas virtualis}, in the
 threefold scheme \latin{circumscriptive}, \latin{definitive}, and
 everywhere (aa.~1--2); two angels cannot be in one place, since two
 complete causes cannot be immediate to one and the same thing (a.~3).}
 {Collect of Compline (\emph{sed contra} of a.~1, as cited); Damascene,
 \emph{De fide orth.} II.3 (\emph{sed contra} and obj.~3 of a.~2);
 Augustine, \emph{De Trin.} VI.6.8; Boethius, \emph{De hebdomadibus}, and
 Aristotle, \emph{Physics}~IV (objections of a.~1); Gen 19:25.}
 {Those who, imagining the angel's indivisibility as a point's, confined
 it to a point-place (a.~2); the objectors who argue from the soul's
 presence in every part of the body (aa.~2--3).}

\subsection{The Local Motion of the Angels (\ST{I q.~53})}\label{sec:q53}

Three articles ask whether an angel can be moved locally, whether in
passing from place to place it passes through the intervening space, and
whether its movement is in time or in an instant (\latin{utrum motus Angeli
sit in tempore vel in instanti}; \ST{I q.~53 pr.}).

\paragraph{Local motion (a.~1).} Aristotle's proof that nothing without
parts is moved is the first objection. The \emph{sed contra} argues from
the Creed: a beatified soul is moved as a beatified angel is, and ``it is
an article of faith that Christ's soul descended into Hell.'' Lateran~IV
confesses that Christ \latin{descendit ad infernos~\dots\ sed descendit in
anima} (DS~801). Aquinas determines that local motion, like being in
place, belongs to body and angel equivocally. Because the angel is in
place ``only by virtual contact,'' its movement ``is nothing else than the
various contacts of various places successively, and not at once.'' The
movement can be continuous, when the angel quits a divisible place part by
part. It can also be non-continuous (\latin{motus non continuus}), when the
angel ``can all at once quit the whole place, and in the same instant apply
himself to the whole of another place.'' Aristotle's demonstration
concerns what is indivisible in quantity and movement that is continuous,
and so does not bind the angel (ad~1). Movement ``by application of
energy is the act of one in act,'' not of an imperfect thing (ad~2).
Nor do the blessed angels move from any need of their own: ``because of
our need, the angel is moved locally,'' according to Heb 1:14, ``Are they
not all ministering spirits, sent to minister for them who shall receive
the inheritance of salvation?'' (ad~3; Douay).

\paragraph{Through the middle or not (a.~2).} The \emph{sed contra} argues
that whatever is changed was first being moved, and so was in the middle.
Aquinas distinguishes the two kinds of angelic motion. If the motion is
continuous, the angel must pass through the middle, since in continuous
motion the order of first and last follows the order of magnitude. If it
is not continuous, the angel can pass ``from one extreme to another without
going through the middle.'' Between any two places lie infinitely many
intermediate places, and only continuous motion traverses such an
infinity; a non-continuous motion would have to number them actually,
which is impossible, so it does not pass through them all. This suits the
angel and not a body, because the angel's substance ``is not subject to
place as contained thereby, but is above it as containing it: hence it is
under his control to apply himself to a place just as he wills, either
through or without the intervening place.'' The second objection compares
the soul, which can think of France and then of Syria without thinking of
Italy. Aquinas replies that the comparison fails, because in local motion
the angel's essence is applied to the places, whereas the things thought
of are in the soul (ad~2).

\paragraph{In time, not in an instant (a.~3).} Some held that the angel's
motion is instantaneous: during the whole preceding time it is in the term
``wherefrom,'' and in the last instant (\latin{instans}) it is in the term
``whereto,'' as with illumination or the generation of fire. Aquinas
answers that to rest is to be in the same disposition in every ``now''
(\latin{nunc}) of a time, so a thing cannot rest in one term throughout a
time and be in the other at its last instant. Instantaneous changes of that
kind are always the terms of some continuous movement, as illumination
ends the local movement of the luminous body. The angel's movement is the
term of no other movement, so ``some `now' must be assigned in which he
was last in the preceding place.'' Where many ``nows'' succeed one another
there is time, ``since time is nothing else than the reckoning of before and
after in movement.'' Hence ``the movement of an angel is in time,''
continuous time if the movement is continuous and non-continuous time if
it is not. Nor is it the time of the heavens that measures all corporeal
things, ``because the angel's movement does not depend upon the movement
of the heavens.'' The replies add two points. First, ``the swiftness of the
angel's movement is not measured by the quantity of his power, but
according to the determination of his will'' (\latin{secundum
determinationem suae voluntatis}; ad~1). Second, in non-continuous time
``an angel can be in one place in one instant, and in another place in the
next instant, without any time intervening'' (ad~3). Illumination is an
alteration terminating a movement, not a local movement, and so is no
model for the angel (ad~2).

\angeldossier{\ST{I q.~53 aa.~1--3}.}
 {Defined: Christ descended to the dead in his soul (Lateran~IV,
 DS~801), the premise of the \emph{sed contra} of a.~1. Common teaching:
 the angels move from place to place in their ministry to men (a.~1 ad~3;
 Heb 1:14). Thomist position: angelic motion is a succession of virtual
 contacts, continuous or non-continuous at the angel's will, able to omit
 the intervening space (aa.~1--2); it is in time, continuous or discrete,
 not in an instant nor in the time of the heavens (a.~3).}
 {Heb 1:14; the article of Christ's descent (DS~801); Aristotle,
 \emph{Physics} III--VI (objections and replies throughout).}
 {Those who held the angel's motion instantaneous, on the model of
 illumination and the generation of fire (a.~3); the Aristotelian
 objection that nothing indivisible can be moved (a.~1). Scotus on the
 ground of the angel's presence in place: \S\ref{sec:dq-place}.}
